\section{Eliminating triangles at equality}
\label{sec:triangles}

The equality witness from Lemma 4.3 has a direct consequence for the
underlying graphon. Its use requires only the positive path-type squares and
first-order vertex variation. We work on a probability space in this section.

\begin{proposition}\label{prop:triangle-elimination}
Let $W$ be a global maximizer of $C(W)/E(W)^{5/2}$ among
graphons of positive edge density, where $C(W)=t_{\mathrm{ind}}(C_5,W)/10$
and $E(W)=t(K_2,W)/2$. If
\[
 p(W)=p_\ast:=\frac5{144},\qquad
 t_{\mathrm{ind}}(C_5,W)=\tau:=\frac5{663552},
\]
then $t(K_3,W)=0$.
\end{proposition}

\begin{proof}
Write $d(x)=\int W(x,y)\,dy$ and
$\kappa=5C(W)/(2E(W))=\tau/(2p_\ast)>0$.
Let $L(x)$ denote one half of the rooted induced-pentagon density at $x$.
Bounded signed variations of the vertex measure, followed by normalization
to total mass one, give
\begin{equation}\tag{32}\label{eq:32}
 L(x)=\kappa d(x)\qquad\text{for almost every }x.
\end{equation}
Indeed, the normalization cancels from the homogeneous ratio, and its
derivative in a bounded direction $a$ is a positive constant times
$\int a(x)(L(x)-\kappa d(x))\,dx$.

At the stated values of $p$ and $t$, the verified certificate is an equality.
Every strictly positively weighted square in its path-type block therefore
vanishes. For ordered root locations $x_0,x_1,x_2$, put
\[
 \rho(x_0,x_1,x_2)=W(x_0,x_1)W(x_1,x_2)\bigl(1-W(x_0,x_2)\bigr),
\]
and let $\phi(x_0,x_1,x_2)$ be the conditional vector of the 72 five-vertex
flags of this rooted type, with the root-type factor omitted. The entries of
$\phi$ are nonnegative and sum to one. If $v_j$ is one of the positive square
vectors, then
\[
 \int\rho\,(v_j\cdot\phi)^2=0.
\]
Thus all $v_j\cdot\phi$ vanish for
$\rho(x_0,x_1,x_2)\,dx_0\,dx_1\,dx_2$-almost every root triple.
Lemma 4.3 supplies a rational linear combination $w$ of these vectors which
is zero on the triangle-free flags and at least one on every flag containing
a triangle. Consequently
\[
 0\leq\sum_{F\text{ containing a triangle}}\phi_F\leq w\cdot\phi=0
\]
on the same set of root triples.

Define
\[
 \begin{split}
 T(x)&=\int W(x,v)W(x,w)W(v,w)\,dv\,dw,\\
 r(x)&=\int W(x,y)W(y,z)\bigl(1-W(x,z)\bigr)\,dy\,dz.
 \end{split}
\]
Among the forbidden flag extensions is the event that the endpoint root
$x_0$ and the two extension vertices span a triangle. Its conditional
probability is exactly $T(x_0)$: edges from the extensions to roots 1 and 2
are unrestricted and are summed out. Hence
\begin{equation}\tag{33}\label{eq:33}
 0=\int\rho(x_0,x_1,x_2)T(x_0)\,dx_0\,dx_1\,dx_2
   =\int r(x)T(x)\,dx.
\end{equation}
In the rooted induced cycle $x-a-b-c-d-x$, retain only the factors
$W(x,a)W(a,b)(1-W(x,b))$. All other factors lie in $[0,1]$, and the two
remaining variables integrate to one. The factor $1/2$ in the definition of
$L$ therefore gives
\[
 L(x)\leq\frac{r(x)}2.
\]
Together with (32), this yields $r(x)\geq2\kappa d(x)$. Since every term
is nonnegative, (33) implies $\int d(x)T(x)\,dx=0$. Finally
$0\leq T(x)\leq d(x)^2$. Thus $T$ vanishes both where $d=0$ and where
$d>0$, up to null sets. It follows that $t(K_3,W)=\int T(x)\,dx=0$.
\end{proof}

\section{Classification of equality}
\label{sec:equality}

The preceding triangle-elimination argument reduces equality in the
unrestricted problem to a triangle-free graphon at the reference density.
We now restrict the positive-square vectors of the same certificate to
triangle-free flags and use the two exact conditional spans recorded in
Section 4.

\subsection{Structure at the reference density}

\begin{proposition}\label{prop:reference-equality}
Let $W$ be a triangle-free global ratio maximizer with
$(p,\theta)=(p_\ast,\tau)$. Then $W$ is a balanced Clebsch graphon on an
active set of measure $1/3$, with isolates outside it.
\end{proposition}

\begin{proof}
By Corollary 4.2, every strictly positive square in the exact certificate
vanishes. For each such square vector $v$, formula (26) gives
$v\cdot q(\mathbf x)=0$ for almost every root tuple with respect to its
type measure.

First use the path type $01,12$. The conditional expectation of the vector
in Lemma 4.4 is $48d(x_1)-5$. Its type probability is
\[
 W(x_0,x_1)W(x_1,x_2)(1-W(x_0,x_2)).
\]
The last factor is redundant almost everywhere on the product of the first
two, by triangle-freeness. Integrating over both endpoints therefore gives
\[
 \int d(x)^2(48d(x)-5)^2\,dx=0.
\]
Thus
\begin{equation}\tag{34}\label{eq:34}
 d(x)\in\left\{0,\frac5{48}\right\}\qquad\text{almost everywhere}.
\end{equation}
Let $A=\{x:d(x)>0\}$. Its complement has zero rows almost everywhere, and
\[
 |A|=\frac{p_\ast}{5/48}=\frac13.
\]
Next consider the type consisting of the edge $01$ and isolated root $2$.
Its two edge endpoints have positive degree for almost every root tuple,
so the conditional expectations of $h_0,h_1$ in Lemma 4.5 vanish by (34).
The same lemma then implies
\begin{equation}\tag{35}\label{eq:35}
 c(x_i,x_2)=\frac25d(x_2),\qquad i=0,1,
\end{equation}
almost everywhere with respect to this type measure. Indeed the conditional
expectation of $k_i$ is $10c(x_i,x_2)-4d(x_2)$.

Restrict to $A$ and normalize its measure to one. Call the resulting
graphon $U$. Its degree is $p_0=5/16$. For type tuples entirely in $A$,
(35) gives codegree $\mu=1/8$ in the new normalization. We now show that
these relations determine $U$.

Write
\[
 c_U(x,z)=\int U(x,y)U(z,y)\,dy.
\]
Triangle-freeness gives $U(x,z)c_U(x,z)=0$ almost everywhere. Integrate
the square of $c_U(x,z)-\mu$ against the edge-plus-isolate type probability
and integrate out its middle root. Since
\[
 \int U(x,y)(1-U(y,z))\,dy=p_0-c_U(x,z),
\]
we obtain
\begin{equation}\tag{36}\label{eq:36}
 (1-U(x,z))(p_0-c_U(x,z))(c_U(x,z)-\mu)^2=0
 \quad\text{almost everywhere}.
\end{equation}
If $0<U(x,z)<1$, triangle-freeness forces $c_U(x,z)=0$, contradicting
(36). Hence $U$ is $\{0,1\}$-valued. Edges have codegree zero, and
nonedges have codegree $\mu$ or $p_0$.

Regard $f_x=U(x,\cdot)$ as an element of $L^2([0,1])$. After a null set
is discarded, all rows are $\{0,1\}$-valued and have integral $p_0$.
Therefore
\[
 \|f_x-f_z\|_2^2=2p_0-2c_U(x,z).
\]
In particular $c_U(x,z)=p_0$ precisely when the two rows agree in $L^2$.
Each twin class $T_x=\{z:f_z=f_x\}$ is measurable: it is the zero set of
the measurable function $z\mapsto\int|U(z,y)-U(x,y)|^2\,dy$.

Let $a(x)=|T_x|$. Almost every twin is a nonneighbor, and every other
nonneighbor has codegree $\mu$. Regularity gives
$\int c_U(x,z)\,dz=p_0^2$, whence
\[
 p_0^2=a(x)p_0+(1-p_0-a(x))\mu.
\]
Substituting $p_0=5/16$ and $\mu=1/8$ gives $a(x)=1/16$ for almost every
$x$. Choose generic representatives successively outside the classes already
selected. Since each selected class has measure $1/16$, exactly sixteen
disjoint measurable classes cover the space up to null sets.

The kernel is constant on every class rectangle. To see this, on a twin
class $A_i$ all rows agree in $L^2$ with one function $f_i$. Fubini and
symmetry give $U(x,z)=f_i(z)=f_j(x)$ almost everywhere on $A_i\times A_j$.
Integrating over either coordinate makes this value constant. Every block
value is zero or one, and diagonal blocks are zero because their pairs have
positive codegree. The quotient is consequently a simple graph on sixteen
vertices, of degree five, with zero common neighbors for adjacent vertices
and two for distinct nonadjacent vertices.

These parameters determine the Clebsch graph elementarily. Fix a vertex
$v$ with independent neighbors $u_1,\ldots,u_5$. Every other vertex has two
neighbors among these five. Each pair $u_i,u_j$ has exactly two common
neighbors, one being $v$; label the other $w_{ij}$. This labels the remaining
ten vertices by two-element subsets of $\{1,\ldots,5\}$. Two pair vertices
sharing a label cannot be adjacent, since that would create a triangle. A
pair vertex needs three further neighbors after $u_i,u_j$, and there are
exactly three disjoint pairs. It joins precisely those three.

Choose the five Clebsch generators $s_1,\ldots,s_5$ from the introduction.
The assignment
\[
 v\mapsto0,\qquad u_i\mapsto s_i,\qquad w_{ij}\mapsto s_i+s_j
\]
is an isomorphism to the Clebsch graph: two pair sums are adjacent exactly
when their pairs are disjoint. Thus $U$ is balanced Clebsch, proving the
proposition.
\end{proof}

\subsection{Arbitrary equality graphons}

\begin{proof}[Equality in Theorem 1.3]
Every member of $\mathcal E_{\mathrm{Cl}}$ attains the bound by the Clebsch
counts and homogeneity. Conversely, suppose $R(W)=R_\ast$ and $E(W)>0$.
The sharp inequality has already been proved for all graphons, so $W$ is
a global maximum over all positive edge densities. Remove its actual
zero-degree set and normalize the remaining vertex measure to one; the
ratio and global maximality are unchanged.

Proposition 3.3 applies to this graphon even if its values are fractional.
Its active density is therefore greater than $p_\ast$. Pad by isolates as
in Lemma 2.5. The padded graphon is still a global maximum, has
$p=p_\ast$ and $\theta=\tau$, and satisfies the vertex first-order
conditions. Proposition 6.1 makes it triangle-free. Proposition 7.1 then
identifies its active part as a balanced Clebsch graphon.

Removing the added isolates and undoing the common vertex-measure scaling
proves $W\in\mathcal E_{\mathrm{Cl}}$. This argument classifies an arbitrary
equality graphon. It does not infer graphon equality from equality on finite
graphs, and it does not impose edge optimality on a merely template-optimal
finite weighting.
\end{proof}

\begin{proof}[Equality in Theorem 1.1]
Normalize the total positive vertex weight to one and represent the weighted
graph as a step graphon with those class measures. At equality, Theorem 1.3
gives a balanced Clebsch active part and isolates. Each original
positive-weight vertex has a constant neighborhood row, so its cell belongs
to one of the twin classes, or to the isolate set, up to null sets. The
sixteen active classes therefore have equal total vertex weight, and their
adjacencies are exactly the Clebsch adjacencies. The converse follows from
the counts $E=40a^2$ and $C=192a^5$ when each active class has total weight
$a$.
\end{proof}

\section{Stability on the edge scale}
\label{sec:stability}

The number of isolated or low-degree vertices is not controlled by an edge
constraint. We first remove a small amount of edge mass to obtain a core
with positive normalized minimum degree. Compactness then places that core
in the neighborhood covered by the local stability theorem. A final
rounding argument produces equal integer class sizes without an additive
error term.

\begin{theorem}[Edge-scaled stability]\label{thm:edge-stability}
There are absolute constants $K<\infty$ and $\varepsilon_0>0$ such that the
following holds. Let $G$ be a finite graph with $m>0$ edges, and suppose that
$0\leq\varepsilon\leq\varepsilon_0$ and
\[
 I_5(G)\geq(R_\ast-\varepsilon)\left(\frac m5\right)^{5/2}.
\]
There is a graph $B$ on the same vertex set, consisting of a Clebsch blow-up
with sixteen classes of the same positive integer size and isolated
vertices, for which
\[
 |E(G)\mathbin{\triangle}E(B)|\leq K\sqrt\varepsilon\,m.
\]
If $s$ is the common class size, then
\[
 |m-40s^2|\leq K\sqrt\varepsilon\,m,\qquad
 s=(1+O(\sqrt\varepsilon))\sqrt{\frac m{40}}.
\]
\end{theorem}

There is no restriction on $|V(G)|$ and no additional rounding term. The
square-root dependence on the deficit is optimal, as shown in Proposition
8.5 below.

\subsection{Removing low-degree vertices}

\begin{lemma}[Weighted peeling]\label{lem:peeling}
Let $(H,x)$ be a finite weighted graph with $E(H,x)=1$ and
$C(H,x)\geq c_\ast-\eta$, where $0\leq\eta\leq c_\ast/4$. Repeatedly remove
a vertex whose current weighted degree is at most $c_\ast$, keeping the
other weights fixed. Write $(H',x')$ for the remaining core,
$E'=E(H',x')$, $C'=C(H',x')$, and
$W'=\sum_{v\in V(H')}x'_v$. Then
\begin{align}
 1-E'&\leq\frac{4\eta}{3c_\ast},\tag{37}\label{eq:37}\\
 \frac{C'}{(E')^{5/2}}&\geq C(H,x),\tag{38}\label{eq:38}\\
 W'&<\frac2{c_\ast}.\tag{39}\label{eq:39}
\end{align}
In particular $E'>0$. After the core weights are normalized to have sum one,
its weighted degrees are all at least $c_\ast^2/2$.
\end{lemma}

\begin{proof}
At an intermediate stage, let $E_t$ be the current edge mass and let $d_i$
be the current weighted degree of a vertex $i$ about to be removed. Fixing
the edge opposite $i$ in a pentagon gives
\[
 C_i\leq\frac14 E_t d_i^2.
\]
Indeed, for a fixed opposite edge $uv$, the two neighbors of $i$ lie in the
exclusive $u$-neighbor and exclusive $v$-neighbor subsets of $N(i)$.
Inducedness makes these sets disjoint even when the graph has triangles.
Their weight product is at most $d_i^2/4$. Removing $i$ destroys edge mass
$x_i d_i$ and pentagon mass $x_i C_i$. As $E_t\leq1$ and $d_i\leq c_\ast$,
the latter is at most $(c_\ast/4)x_i d_i$.

Put $\alpha=1-E'$. Each removed edge and each destroyed pentagon is charged
at its first deleted endpoint, so
\begin{equation}\tag{40}\label{eq:40}
 0\leq C(H,x)-C'\leq\frac{c_\ast}4\alpha.
\end{equation}
Theorem 1.1, applied to the surviving core, gives
\[
 c_\ast-\eta-\frac{c_\ast}4\alpha\leq C'
 \leq c_\ast(1-\alpha)^{5/2}.
\]
Since $1-(1-\alpha)^{5/2}\geq\alpha$ for $0\leq\alpha\leq1$, this proves
(37), and hence $E'\geq2/3$. Moreover $C(H,x)\geq3c_\ast/4$, so
\[
 C'\geq C(H,x)-\frac{c_\ast}4\alpha
 \geq C(H,x)(1-\alpha)\geq C(H,x)(1-\alpha)^{5/2}.
\]
This is (38).

Every surviving vertex has weighted degree greater than $c_\ast$.
Consequently
\[
 c_\ast W'<\sum_{i\in V(H')}x_i d_i=2E'\leq2.
\]
After dividing the vertex weights by $W'$, every weighted degree exceeds
\[
 \frac{c_\ast}{W'}>\frac{c_\ast^2}{2E'}\geq\frac{c_\ast^2}2.
\]
The degrees must be recomputed after each removal; this is why the procedure
is iterative.
\end{proof}

\subsection{Entering the local neighborhood}
Let $W_{\mathrm{Cl}}$ denote a full balanced Clebsch graphon. We use
$\delta_\square$ for cut distance and $\delta_1$ for $L^1$ distance, in each
case taking the infimum over measure-preserving relabelings. The standard
compactness of graphons in cut distance and continuity of fixed finite
homomorphism densities will be used only qualitatively.

\begin{lemma}\label{lem:positive-degree-compactness}
Fix $\mu>0$. If $W_j$ are graphons with $d_{W_j}\geq\mu$ almost everywhere
and $R(W_j)\to R_\ast$, then $\delta_1(W_j,W_{\mathrm{Cl}})\to0$.
\end{lemma}

\begin{proof}
Their edge densities are at least $\mu$. From any subsequence, compactness
extracts a further subsequence converging in cut distance to a graphon $W$.
Choose relabelings giving cut-norm convergence. For every measurable
$S\subseteq[0,1]$,
\[
 \int_S d_W=\lim_j\int_S d_{W_j}\geq\mu|S|,
\]
because the difference of the two integrals is bounded by the cut norm.
Thus $d_W\geq\mu$ almost everywhere. The equality $R(W)=R_\ast$ passes to
the limit by induced-density continuity; the denominator in $R$ stays
bounded away from zero. Theorem 1.3 identifies $W$ as balanced Clebsch on
its active set, with isolates outside. The minimum-degree condition
excludes a positive-measure isolate set, so $W$ is a full balanced Clebsch
graphon.

This proves convergence of the whole sequence in cut distance. To obtain
$L^1$ convergence, fix the sixteen classes $A_1,\ldots,A_{16}$ of a labeling
of $W_{\mathrm{Cl}}$. Since $W_{\mathrm{Cl}}$ is $\{0,1\}$-valued,
$U-W_{\mathrm{Cl}}$ has a constant sign on each $A_i\times A_j$. Hence
\[
 \|U-W_{\mathrm{Cl}}\|_1
 =\sum_{i,j}\left|\int_{A_i\times A_j}(U-W_{\mathrm{Cl}})\right|
 \leq256\|U-W_{\mathrm{Cl}}\|_\square.
\]
Taking infima over relabelings completes the proof.
\end{proof}

\subsection{Rounding whole vertices}
Write $\mathcal E_{\mathrm{Cl}}$ for the family defined in the introduction.
The next lemma is useful because balancing real class masses by rounding
their sizes separately would introduce an unnecessary additive error.

\begin{lemma}[Integral rounding]\label{lem:integral-rounding}
Let $W_G$ be the probability graphon of an $n$-vertex graph, and let
$U\in\mathcal E_{\mathrm{Cl}}$ have active mass $q\geq1/2$. There is a graph
$B$ on $V(G)$, consisting of sixteen equal integer Clebsch classes and
isolates, such that
\[
 \|W_G-W_B\|_1\leq897\|W_G-U\|_1.
\]
The common class size is positive when $\|W_G-U\|_1$ is sufficiently small.
\end{lemma}

\begin{proof}
Put $\zeta=\|W_G-U\|_1$. Let $B_0$ be the isolate set of $U$, and let
$B_1,\ldots,B_{16}$ be its active classes, each of measure $b=q/16$.
Denote the corresponding row functions by $r_0=0,r_1,\ldots,r_{16}$.
Different rows have distance at least
\begin{equation}\tag{41}\label{eq:41}
 \rho=\frac{5q}{16}\geq\frac5{32}.
\end{equation}
A nonzero row meets five classes, while distinct Clebsch rows differ on six
or ten classes.

For a vertex $v$, let $I_v$ be its cell in $W_G$, of measure $1/n$, and let
$g_v$ be its constant row. Set
\[
 p_{v,i}=n|I_v\cap B_i|,\qquad
 e_{v,i}=\|g_v-r_i\|_1,\qquad
 e_v=\sum_{i=0}^{16}p_{v,i}e_{v,i}.
\]
Assign the entire vertex $v$ to a type $i(v)$ minimizing $e_{v,i}$. If this
minimum is at least $\rho/2$, then $e_v\geq\rho/2$. Otherwise, the triangle
inequality and (41) give $e_{v,j}\geq\rho/2$ for every $j\ne i(v)$. In both
cases
\[
 1-p_{v,i(v)}\leq\frac{2e_v}{\rho}.
\]
The total measure whose type changes is therefore
\begin{equation}\tag{42}\label{eq:42}
 M=\frac1n\sum_v(1-p_{v,i(v)})\leq\frac{2\zeta}{\rho}.
\end{equation}
Let $T$ be the Clebsch step graphon given by these whole-vertex assignments,
and let $s_i$ be the mass of its $i$th active class. Changing types on mass
$M$ gives
\[
 \|T-U\|_1\leq2M,\qquad\sum_{i=1}^{16}|s_i-b|\leq2M.
\]
Set $s_{\min}=\min_i s_i$. Remove excess vertices from each active class
until every class has size $ns_{\min}$, assigning the removed vertices to
the isolate type. This uses whole vertices because all $ns_i$ are integers.
The total mass reassigned is
\[
 \Delta=\sum_{i=1}^{16}(s_i-s_{\min})
 \leq17\sum_{i=1}^{16}|s_i-b|\leq34M.
\]
For the first inequality, if $b\geq s_{\min}$ use
$b-s_{\min}\leq\sum_i|s_i-b|$; if $b<s_{\min}$, the inequality is immediate.
The additional $L^1$ cost is at most $2\Delta\leq68M$. The resulting graph
$B$ is exactly balanced, and
\[
 \|W_G-W_B\|_1\leq\zeta+2M+68M
 \leq\left(1+\frac{140}{\rho}\right)\zeta\leq897\zeta.
\]
For sufficiently small $\zeta$, (42) and $b\geq1/32$ also imply $s_i>0$
for every $i$.
\end{proof}

This proof applies after arbitrary relabelings: the cells $I_v$ need not be
intervals. It also handles a nonattained distance infimum. Apply the lemma
to a sequence of comparison graphons approaching the infimum and use the
finiteness of the possible edit distances on a fixed vertex set. If the
infimum is zero, taking the error small enough to permit fewer than one
edge edit forces an exact copy.

\subsection{Proof of the stability theorem}

\begin{proof}[Proof of Theorem 8.1]
Give every vertex of $G$ weight $1/\sqrt m$. Then $E=1$ and
\[
 C\geq c_\ast-\eta,\qquad\eta=\frac{\varepsilon}{5^{5/2}}.
\]
Apply Lemma 8.2. The peeling removes at most
\begin{equation}\tag{43}\label{eq:43}
 \frac{4\varepsilon}{3c_\ast5^{5/2}}m
\end{equation}
edges and leaves a graph $H$ with
\[
 n_H<\frac{2\sqrt m}{c_\ast}.
\]
Its probability graphon has minimum degree at least $c_\ast^2/2$ and ratio
at least $R_\ast-\varepsilon$.

By Lemma 8.3, there is an absolute positive threshold for $\varepsilon$
below which $W_H$, after a relabeling, lies in the neighborhood of Theorem
A.1. Let $c_{\mathrm{loc}}>0$ be the constant in that theorem. Its
induced-pentagon functional is exactly the functional used here, so the
local conclusion is
\[
 d_1(W_H,\mathcal E_{\mathrm{Cl}})\leq
 \sqrt{\frac{\varepsilon}{c_{\mathrm{loc}}}}.
\]
Comparison graphons sufficiently close to this infimum have active mass
at least $1/2$ when $\varepsilon$ is small. To see this, also choose the
compactness neighborhood small enough that $W_H$ is close to the full
Clebsch graphon in $L^1$. Edge-density continuity and
$t(K_2,U)=(5/16)q^2$ then force $q$ close to one.

Lemma 8.4 gives a balanced Clebsch blow-up $B_H$ with isolates on $V(H)$.
Taking the comparison distance to its infimum, the number of edits inside
$H$ is at most
\[
 \frac{897n_H^2}2\sqrt{\frac{\varepsilon}{c_{\mathrm{loc}}}}
 \leq\frac{1794}{c_\ast^2\sqrt{c_{\mathrm{loc}}}}
 \sqrt\varepsilon\,m.
\]
Extend $B_H$ to the original vertex set by making every peeled vertex
isolated. Adding (43), and using $\varepsilon\leq\sqrt\varepsilon$ for
$\varepsilon\leq1$, proves the edit bound. For example, once
$\varepsilon_0$ has been chosen small enough for the preceding neighborhood
and nonemptiness requirements, one may take
\[
 K=\frac{1794}{c_\ast^2\sqrt{c_{\mathrm{loc}}}}
   +\frac4{3c_\ast5^{5/2}}.
\]
The value of $K$ is not optimized. Finally, a balanced Clebsch blow-up with
common class size $s$ has $40s^2$ edges. Its edge count differs from $m$ by
at most the number of edits, proving both remaining conclusions.
\end{proof}

\subsection{Optimality of the exponent}

\begin{proposition}\label{prop:sqrt-optimality}
The factor $\sqrt\varepsilon$ in Theorem 8.1 cannot be replaced by
$o(\sqrt\varepsilon)$, uniformly over finite graphs.
\end{proposition}

\begin{proof}
Choose two nonadjacent vertices $u,v$ of $\mathrm{Cl}$. For integers
$k\geq2$ and $1\leq r<k$, form a Clebsch blow-up $G_{k,r}$ with class sizes
$k+r$ at $u$, $k-r$ at $v$, and $k$ at the other fourteen vertices.

Each Clebsch vertex lies in sixty pentagons, and each nonadjacent pair lies
in twelve. For the latter count, the pair has two common neighbors and six
three-edge paths between its endpoints. Each choice of a common neighbor
and such a path forms a pentagon, by triangle-freeness. The six paths can
also be counted from the five neighbors of one endpoint: two are adjacent
to the other endpoint and contribute zero codegree, while the other three
contribute codegree two.

The edge and pentagon polynomials are affine in each class size. The linear
changes at $u$ and $v$ cancel, and there is no quadratic edge term because
$u,v$ are nonadjacent. Thus, exactly,
\[
 |V(G_{k,r})|=16k,\qquad m=40k^2,\qquad
 I_5(G_{k,r})=192k^5-12k^3r^2.
\]
Its normalized deficit is
\[
 \varepsilon_{k,r}=R_\ast-\frac{I_5(G_{k,r})}{(m/5)^{5/2}}
 =\frac{R_\ast r^2}{16k^2}.
\]
The pair $u,v$ has two common neighbors. Its three exclusive neighbors on
each side have unchanged class size $k$, so $G_{k,r}$ has $3k$ vertices of
degree $5k+r$, $3k$ of degree $5k-r$, and $10k$ of degree $5k$.

A balanced Clebsch blow-up with isolates on the same $16k$ vertices has
common class size at most $k$, and hence maximum degree at most $5k$.
Each of the $3k$ high-degree vertices must lose at least $r$ in degree. One
edge edit changes the total absolute degree discrepancy by at most two.
Therefore every such target $B$ satisfies
\[
 |E(G_{k,r})\mathbin{\triangle}E(B)|\geq\frac32kr
 =\frac3{20\sqrt{R_\ast}}\sqrt{\varepsilon_{k,r}}\,m.
\]
Taking $r=1$ and $k\to\infty$ proves the assertion.
\end{proof}

\section{Linear structural stability and quadratic mass stability}
\label{sec:linear-stability}

The local repair argument proves a stronger conclusion than its quadratic
distance estimate alone. Edge errors have a linear cost in the extremal
deficit; only the imbalance of the sixteen class masses has a quadratic
cost. We retain these two conclusions separately.

Let $\mathcal T_{\mathrm{Cl}}$ denote the class of graphons obtained from an
arbitrary sixteen-class Clebsch blow-up and an isolated class. The sixteen
active masses need not be equal. For a presented member $T$ with active
masses $s_1,\ldots,s_{16}$, put
\[
 q(T)=\sum_{i=1}^{16}s_i,\qquad
 V(T)^2=\sum_{i=1}^{16}(s_i-q(T)/16)^2.
\]

\begin{proposition}[The two local defects]\label{prop:two-defects}
There are constants $A,B>0$ and $\delta>0$ such that, whenever
$\|W-W_{\mathrm{Cl}}\|_1<\delta$, a graphon $T\in\mathcal T_{\mathrm{Cl}}$
exists with all active masses at least $1/32$ and
\begin{equation}\tag{44}\label{eq:44}
 \|W-T\|_1\leq AD(W),\qquad V(T)^2\leq BD(W),\qquad
 D(W)=R_\ast-R(W).
\end{equation}
\end{proposition}

\begin{proof}
Use the sets $B_0,Z,M_1,\ldots,M_{16}$ constructed in the proof of Theorem
A.1; here $B_0$ denotes its bad-row set to avoid confusion with the constant
in the present statement. Write $F=C/E^{5/2}=R/5^{5/2}$ as in that proof.
Take $T$ to be the repaired Clebsch step graphon on $M=\bigcup_i M_i$, and
put isolates outside $M$. Before balancing the masses, the four-defect
inequality is
\[
 F(W_{\mathrm{Cl}})-F(W)
 \geq c_B\mu_{B_0}+c_ZD_Z+c_Ee_M+c_HV(T)^2.
\]
The three operations that produce $T$ cost at most
$2\mu_{B_0}+2D_Z+e_M$ in $L^1$. Consequently one may take
\[
 A=\frac1{5^{5/2}\min\{c_B/2,c_Z/2,c_E\}},\qquad
 B=\frac1{5^{5/2}c_H}.
\]
The mass-closeness estimate in that proof gives
$s_i=1/16+O(\|W-W_{\mathrm{Cl}}\|_1)$ uniformly in $i$. Decreasing the
fixed neighborhood ensures $s_i\geq1/32$.
\end{proof}

\begin{lemma}[Rounding without balancing]\label{lem:unbalanced-rounding}
Let $W_G$ be the probability graphon of a graph on $n$ vertices, and let
$U\in\mathcal T_{\mathrm{Cl}}$ have active masses $s_i\geq1/32$. Set
$\zeta=\|W_G-U\|_1$ and $\rho=5/32$. There is a Clebsch blow-up $H$, with
isolates, on the same vertex set, such that
\begin{align}
 \|W_G-W_H\|_1&\leq(1+4/\rho)\zeta,\tag{45}\label{eq:45}\\
 V(W_H)^2&\leq2V(U)^2+32\zeta^2/\rho^2.\tag{46}\label{eq:46}
\end{align}
All sixteen active classes of $H$ are nonempty if $\zeta<\rho/128$.
\end{lemma}

\begin{proof}
The seventeen row prototypes of $U$ are pairwise separated in $L^1$ by at
least $\rho$. Indeed, a nonzero Clebsch row meets five classes, while distinct
nonzero rows differ on at least six classes.

For each vertex cell $I_v$, write $p_{v,i}=n|I_v\cap A_i|$ and
$e_{v,i}=\|W_G(v,\cdot)-r_i\|_1$, including the isolated type $i=0$.
Assign the whole vertex to a type minimizing $e_{v,i}$. The argument of
Lemma 8.4 gives total changed type mass
\[
 M=\frac1n\sum_v(1-p_{v,i(v)})\leq\frac{2\zeta}{\rho}.
\]
The induced whole-vertex template $W_H$ differs from $U$ only on pairs
with at least one endpoint in this changed set. Therefore
$\|U-W_H\|_1\leq2M$, proving (45). No trimming is performed.

Let $s'$ be the vector of the sixteen rounded active masses. The changed
mass estimate implies $\|s'-s\|_1\leq2M$. Let $P$ be orthogonal projection
onto the sum-zero subspace of $\mathbb R^{16}$. Then
\[
 \|Ps'\|_2^2\leq2\|Ps\|_2^2+2\|P(s'-s)\|_2^2
 \leq2V(U)^2+8M^2,
\]
which is (46). Finally
$s'_i\geq s_i-2M\geq1/32-4\zeta/\rho>0$ under the stated smallness
assumption. Each $ns'_i$ is an integer because only whole vertices were
assigned. The vertex cells can be arbitrary measurable sets, so the argument
also applies after a graphon relabeling.
\end{proof}

\begin{theorem}[Linear structural stability]\label{thm:linear-stability}
There are absolute constants $K<\infty$ and $\varepsilon_0>0$ such that the
following holds. Suppose $G$ has $m>0$ edges and
\[
 I_5(G)\geq(R_\ast-\varepsilon)(m/5)^{5/2},\qquad
 0\leq\varepsilon\leq\varepsilon_0.
\]
Then a graph $H$ on the same vertex set, consisting of a Clebsch blow-up with
positive integer class sizes $n_1,\ldots,n_{16}$ and isolated vertices,
satisfies
\begin{equation}\tag{47}\label{eq:47}
 |E(G)\mathbin{\triangle}E(H)|+\sum_{i=1}^{16}(n_i-N/16)^2
 \leq K\varepsilon m,\qquad N=\sum_{i=1}^{16}n_i.
\end{equation}
In particular,
\begin{equation}\tag{48}\label{eq:48}
 \left|m-\frac{5N^2}{32}\right|\leq K'\varepsilon m,\qquad
 N=\sqrt{\frac{32m}5}+O(\varepsilon\sqrt m),
\end{equation}
and every class satisfies
\[
 n_i=(1+O(\sqrt\varepsilon))\sqrt{m/40}.
\]
All constants are absolute. There is no additive rounding term.
\end{theorem}

\begin{proof}
Give every vertex weight $1/\sqrt m$. The weighted peeling lemma removes
at most $4\varepsilon m/(3R_\ast)$ edges, and leaves a core $G'$ with
$n<2\sqrt m/c_\ast$ vertices, normalized minimum degree at least
$c_\ast^2/2$, and ratio at least $R_\ast-\varepsilon$. By the positive-degree
compactness lemma, after relabeling its probability graphon lies in the
neighborhood of Proposition 9.1 when $\varepsilon_0$ is small enough.
Obtain the template $U$ there. It satisfies
\[
 \zeta\leq A\varepsilon,\qquad V(U)^2\leq B\varepsilon.
\]
Decrease $\varepsilon_0$ further so that $A\varepsilon_0<\rho/128$ and
$\varepsilon_0\leq1$. Lemma 9.2 gives whole-vertex classes with
\[
 \begin{split}
 \|W_{G'}-W_H\|_1&\leq(1+4/\rho)A\varepsilon,\\
 V(W_H)^2&\leq B'\varepsilon,\qquad B'=2B+32A^2/\rho^2.
 \end{split}
\]
The number of edge edits inside the core is $n^2/2$ times its $L^1$
distance. Also $\sum_i(n_i-N/16)^2=n^2V(W_H)^2$. Using
$n^2<4m/c_\ast^2$ and putting all peeled vertices into the isolated class
proves the two estimates separately with constants
\[
 K_E=\frac4{3R_\ast}+\frac{2(1+4/\rho)A}{c_\ast^2},\qquad
 K_V=\frac{4B'}{c_\ast^2}.
\]
Take $K=K_E+K_V$ in (47).

For completeness, let $A_{\mathrm{Cl}}$ be the Clebsch adjacency matrix
and $h_i=n_i-N/16$. Since the graph is five-regular and $\sum_i h_i=0$,
\[
 e(H)=40(N/16)^2+\frac12 h^TA_{\mathrm{Cl}}h
 =\frac{5N^2}{32}+\frac12 h^TA_{\mathrm{Cl}}h.
\]
The elementary bound $|h^TA_{\mathrm{Cl}}h|\leq5\|h\|_2^2$ follows from
$2|h_i h_j|\leq h_i^2+h_j^2$ summed over the edges. Hence one may take
$K'=K_E+(5/2)K_V$ in the first estimate of (48). For smaller
$\varepsilon_0$, this places $N$ between fixed positive multiples of
$\sqrt m$. Factoring the difference of squares yields the second estimate.
Finally $|h_i|\leq\|h\|_2=O(\sqrt{\varepsilon m})$, which gives the
class-size conclusion.
\end{proof}

\begin{proposition}[Optimality of linear structural stability]
\label{prop:linear-optimality}
The factor $\varepsilon$ in the edit-distance conclusion of Theorem 9.3
cannot be replaced uniformly by $o(\varepsilon)$.
\end{proposition}

\begin{proof}
Start with a balanced Clebsch blow-up of common class size $k\geq2$, and
add a single edge between two vertices in one class. The two endpoints are
now true twins. An induced pentagon cannot contain both, and all other
induced pentagons are unchanged. Thus the new graph has
\[
 m_k=40k^2+1,\qquad I_5=192k^5,\qquad
 \varepsilon_k=R_\ast\left[1-\left(\frac{40k^2}{40k^2+1}\right)^{5/2}\right].
\]
Its distance to the family of Clebsch blow-ups and isolates is exactly one:
deleting the added edge suffices, and the graph contains a triangle whereas
every graph in that family is triangle-free. Furthermore $\varepsilon_k\to0$
and $\varepsilon_km_k\to(5/2)R_\ast$. A bound of the form
$o(\varepsilon_k)m_k$ would tend to zero, contradicting this integer
distance of one.
\end{proof}

\subsection{A two-sided geometric description of the deficit}
For a graph $G$ with $m>0$ edges, write
\[
 \Delta(G)=c_\ast m^{5/2}-I_5(G).
\]
Let $\mathcal B_m(G)$ consist of the Clebsch blow-ups with isolates on
$V(G)$ whose sixteen class sizes obey
\[
 \frac12\sqrt{m/40}\leq n_i\leq2\sqrt{m/40}.
\]
For such a template put $N=\sum_i n_i$ and define
\[
 J_G(B)=|E(G)\mathbin{\triangle}E(B)|+\sum_i(n_i-N/16)^2.
\]

\begin{proposition}[Two-sided structural energy]\label{prop:structural-energy}
For absolute constants $c,C>0$, every graph sufficiently close to the sharp
ratio satisfies
\begin{equation}\tag{49}\label{eq:49}
 c\,m^{3/2}\min_{B\in\mathcal B_m(G)}J_G(B)
 \leq\Delta(G)\leq
 C\,m^{3/2}\min_{B\in\mathcal B_m(G)}J_G(B).
\end{equation}
In particular $\mathcal B_m(G)$ is nonempty in this range. The upper
estimate holds whenever $\mathcal B_m(G)$ is nonempty, without a separate
near-extremality assumption.
\end{proposition}

\begin{proof}
For a sufficiently small fixed normalized deficit, the class-size conclusion
of Theorem 9.3 places its template in $\mathcal B_m(G)$. Substituting
$\varepsilon=5^{5/2}\Delta(G)/m^{5/2}$ in that theorem proves the lower
estimate in (49).

For the other direction, fix $B\in\mathcal B_m(G)$ and put
\[
 t=|E(G)\mathbin{\triangle}E(B)|,\qquad
 V=\sum_i(n_i-N/16)^2,\qquad m_B=e(B).
\]
An original induced pentagon of $B$ is still induced in $G$ unless one of
its ten pairs is modified. Each modified pair belongs to at most
$\binom{N-2}3\leq N^3/6$ original five-sets; pairs meeting an isolated
vertex of $B$ belong to none. Thus
\begin{equation}\tag{50}\label{eq:50}
 I_5(G)\geq I_5(B)-tN^3/6.
\end{equation}
This estimate depends on the number of active template vertices, not on the
possibly much larger number of vertices of $G$.

Consider the weighted-template defect
\[
 Q(x)=c_\ast E(x)^{5/2}-P(x).
\]
At every positive constant vector $b\mathbf1$ we have $Q=0$ and
$\nabla Q=0$. The latter follows either by differentiating the sharp equality
along class weights, or directly from $E(\mathbf1)=40$, $P(\mathbf1)=192$,
$\nabla E(\mathbf1)=5\mathbf1$ and $\nabla P(\mathbf1)=60\mathbf1$.
Set $b=N/16$ and $h_i=n_i-b$. Every point of the line segment
$b\mathbf1+uh$, $0\leq u\leq1$, has coordinates between fixed positive
multiples of $\sqrt m$. On this segment the Hessian of $Q$ is
$O(m^{3/2})$: the polynomial $P$ has degree five, and the same scaling
follows directly by differentiating $E^{5/2}$, with $E$ bounded above and
below by fixed positive multiples of $m$. Taylor's formula therefore gives
\[
 0\leq c_\ast m_B^{5/2}-I_5(B)=Q(b\mathbf1+h)
 \leq C_0m^{3/2}\|h\|_2^2=C_0m^{3/2}V.
\]
Also $m_B$ and $m$ are within fixed positive constant factors, and
$|m_B-m|\leq t$, so
\[
 |m^{5/2}-m_B^{5/2}|\leq C_1m^{3/2}t.
\]
Since $N=O(\sqrt m)$, combine these estimates with (50) to obtain
$\Delta(G)\leq Cm^{3/2}(t+V)$. Taking the minimum completes the proof.
There are finitely many vertex partitions, so the minimum is attained
whenever the family is nonempty.
\end{proof}

\section{The arithmetic size of the finite-edge deficit}
\label{sec:arithmetic-size}

The error term in Corollary 1.2 cannot be improved uniformly to $o(m^2)$.
More precisely, its order is governed by the distance of an effective vertex
count from the integers. This conclusion uses the linear stability theorem
for unbalanced Clebsch blow-ups, including its squared class-size estimate.

For an integer $N\geq0$, write $N=16k+r$, with $0\leq r<16$, and put
\[
 v(N)=\frac{r(16-r)}{16},\qquad
 \Psi(m)=\min_{N\in\mathbb Z_{\geq0}}
 \left(\left|m-\frac5{32}N^2\right|+v(N)\right).
\]
The minimum is attained, and $0\leq v(N)\leq4$. As before, let
$\rho(C_5,m)$ denote the maximum number of induced pentagons among graphs
with exactly $m$ edges, and set
\[
 \Delta(m)=c_\ast m^{5/2}-\rho(C_5,m).
\]

\begin{theorem}[Arithmetic deficit law]\label{thm:arithmetic-deficit}
There exist absolute constants $a_0,b_0>0$ and an integer $m_0$ such that,
for every integer $m\geq m_0$,
\begin{equation}\tag{51}\label{eq:51}
 a_0m^{3/2}\Psi(m)\leq\Delta(m)\leq b_0m^{3/2}\Psi(m).
\end{equation}
In particular, $\Psi(m)=0$ exactly when $m=40k^2$ for an integer $k\geq0$.
For all other sufficiently large integers $m$,
\begin{equation}\tag{52}\label{eq:52}
 \Delta(m)\asymp m^2\operatorname{dist}
 \left(\sqrt{\frac{32m}5},\mathbb Z\right)+m^{3/2},
\end{equation}
where the two comparison constants are absolute.
\end{theorem}

\begin{proof}
We first record the role of the integer class sizes. If
$n_1,\ldots,n_{16}$ are nonnegative integers with sum $N$, then
\begin{equation}\tag{53}\label{eq:53}
 V:=\sum_{i=1}^{16}(n_i-N/16)^2\geq v(N).
\end{equation}
Indeed, whenever two entries differ by at least two, moving one unit from
the larger to the smaller decreases the sum of their squares. The minimum
is therefore attained by $r$ entries equal to $k+1$ and $16-r$ entries equal
to $k$; its value is $v(N)$.

Let $A$ be the adjacency matrix of $\mathrm{Cl}$, and write
$n=(N/16)\mathbf1+h$, where $\sum_i h_i=0$. Since $A\mathbf1=5\mathbf1$,
the edge count of this blow-up is
\begin{equation}\tag{54}\label{eq:54}
 E(n)=\frac12n^TAn=\frac5{32}N^2+\frac12h^TAh.
\end{equation}
The operator norm of $A$ is at most five. Consequently
\begin{equation}\tag{55}\label{eq:55}
 \left|E(n)-\frac5{32}N^2\right|\leq\frac52 V.
\end{equation}
For the lower bound in (51), choose an extremal graph $G$ with $m$ edges.
Such a graph exists, since discarding isolates leaves at most $2m$ vertices.
The earlier lower construction gives $\Delta(m)=O(m^2)$. Hence its
normalized deficit
\[
 \varepsilon=R_\ast-\frac{I_5(G)}{(m/5)^{5/2}}
 =\frac{5^{5/2}\Delta(m)}{m^{5/2}}
\]
tends to zero as $m\to\infty$. Apply Theorem 9.3. It supplies a Clebsch
blow-up $H$, together with isolates, on the same vertex set, such that its
integer class sizes $n_i$, total active size $N$, squared deviation $V$, and
edge-edit distance $L$ satisfy
\[
 L+V\leq K\varepsilon m=K5^{5/2}\frac{\Delta(m)}{m^{3/2}}.
\]
Here and below a constant may be increased without changing its dependence.
In particular $|m-E(n)|\leq L$. Equations (53) and (55) give
\[
 \Psi(m)\leq\left|m-\frac5{32}N^2\right|+v(N)
 \leq L+\frac72V\leq K'\frac{\Delta(m)}{m^{3/2}}.
\]
This proves the first inequality.

For the upper bound, put $t=\sqrt{32m/5}$. An integer $N_0$ nearest to $t$
satisfies
\[
 \left|m-\frac5{32}N_0^2\right|
 \leq\frac5{32}\left(t+\frac14\right).
\]
It follows that $\Psi(m)=O(\sqrt m)$. Any minimizing $N$ therefore
satisfies $N=\Theta(\sqrt m)$ and tends to infinity. Choose
$n_1,\ldots,n_{16}$ as equally as possible, so that $V=v(N)$, and let $H$
be the resulting Clebsch blow-up. With $e=E(n)$, (55) gives
\begin{equation}\tag{56}\label{eq:56}
 |e-m|\leq\left|m-\frac5{32}N^2\right|+\frac52v(N)
 \leq\frac52\Psi(m).
\end{equation}
We also need the second-order cost of the integer balancing. For positive
real class sizes $x$, define
\[
 Q(x)=c_\ast E(x)^{5/2}-P(x).
\]
The Clebsch counts give $Q(\mathbf1)=0$ and $\nabla Q(\mathbf1)=0$. The
function $Q$ has a bounded Hessian in a fixed neighborhood of $\mathbf1$.
Write $q=N/16$ and $n=q\mathbf1+h$, where $|h_i|\leq1$ and
$\|h\|_2^2=v(N)$. Homogeneity and Taylor's theorem yield, uniformly for
large $N$,
\begin{equation}\tag{57}\label{eq:57}
 0\leq Q(n)=q^5Q(\mathbf1+h/q)
 \leq Cq^3\|h\|_2^2\leq C'N^3v(N).
\end{equation}
The nonnegativity follows from the sharp weighted bound. If $v(N)=0$, the
same equation holds exactly, with both sides zero.

If $e\leq m$, add a disjoint matching with $m-e$ edges. If $e>m$, delete
any $e-m$ edges of $H$. A single edge deletion can destroy at most
$\binom{N-2}3$ induced pentagons, since every affected pentagon must contain
both endpoints of that edge. Thus the resulting graph $G_m$ has exactly
$m$ edges and
\[
 I_5(G_m)\geq P(n)-|e-m|\binom{N-2}3.
\]
Since $e=m+O(\sqrt m)$ and $N=\Theta(\sqrt m)$, the mean value theorem,
(56), and (57) imply
\[
 \begin{split}
 \Delta(m)&\leq c_\ast m^{5/2}-I_5(G_m)\\
 &\leq c_\ast|m^{5/2}-e^{5/2}|+Q(n)+|e-m|\binom{N-2}3\\
 &\leq C''m^{3/2}\bigl(|e-m|+v(N)\bigr)
 \leq b_0m^{3/2}\Psi(m).
 \end{split}
\]
This proves (51).

Finally, $\Psi(m)=0$ requires $v(N)=0$, hence $16\mid N$, and then
$m=5N^2/32=40(N/16)^2$. The converse is immediate. If $m$ is not of this
form, every term in the minimum is at least $15/16$: for $16\mid N$ the
nonzero integer $|m-5N^2/32|$ is at least one, and otherwise
$v(N)\geq15/16$. Let $d=\operatorname{dist}(t,\mathbb Z)$. For every
$N\geq0$,
\[
 \left|m-\frac5{32}N^2\right|=\frac5{32}|t-N|(t+N)
 \geq\frac5{32}td.
\]
Choosing an integer nearest to $t$ gives
\[
 \Psi(m)\leq\frac5{32}(2t+1/2)d+4.
\]
Together with the lower bound $15/16$, these inequalities show
$\Psi(m)\asymp\sqrt m\,d+1$ outside the equality sequence. Equation (52)
follows.
\end{proof}

\begin{corollary}[The $m^2$ error is uniformly optimal]
\label{cor:arithmetic-optimality}
For $m_k=40k^2+\lfloor5k/2\rfloor$,
\[
 \Delta(m_k)=\Theta(m_k^2).
\]
For each fixed nonzero integer $j$ and sufficiently large $k$,
\[
 \Delta(40k^2+j)=\Theta\bigl(|j|(40k^2+j)^{3/2}\bigr).
\]
At $m=40k^2$ the deficit is exactly zero.
\end{corollary}

\begin{proof}
For the first sequence, the minimizing index in $\Psi$ is $N=16k$, and
$\Psi(m_k)=\lfloor5k/2\rfloor$. Indeed, the next larger square value is
$40k^2+5k+5/32$, and all smaller indices have a strictly greater distance
from $m_k$ than $N=16k$; the nonnegative term $v(N)$ cannot reverse this.
For fixed $j\ne0$, the same index minimizes for all sufficiently large
$k$, with $\Psi(40k^2+j)=|j|$. Apply Theorem 10.1. The last assertion is
the exact Clebsch construction.
\end{proof}

\begin{remark}\label{rem:arithmetic-scale}
Theorem 10.1 determines the order of the finite-edge deficit, with constants
uniform in $m$. It does not identify the exact maximizing graph at each
integer edge count. Theorem 12.2 below identifies the leading coefficient
on the $m^2$ scale. Its uniform absolute remainder does not replace the
present comparison at phases extremely close to zero.
\end{remark}
